\documentstyle[% 


righttag% 


,11pt% 


,amssymb% 


,russian%               remove in Eng. 


,izvru%                 Set izven% in Eng. 


]{amsart} 


 


%  


 


\newcommand{\tts}[1]{} 


 


\theoremstyle{plain} 


\theorembodyfont{\it} 


\newtheorem{thmnonum}{\hskip\parindent Теорема} 


\renewcommand{\thethmnonum}{} 


 


 


%\hoffset=-15mm%                  For Eng. 


\setcounter{page}{81} 


\god{2004} 


\nomer{3~(502)} %                        Remove in Eng. 


%% remove ^^^^^^ 


%\nomer{Vol.\,48, No.\,3}%               Set in Eng. 


%\str{\pageref{begin}--\pageref{end}}%  Set in Eng. 


\udk{519.71} 


 


\author{\it И.К.\,Шаранхаев} 


% NB:   ^^ remove in Eng. 


\title{О булевых базисах второго яруса} 


 


\date{} 


%\pagestyle{myheadings}%         Set in Eng. 


\pagestyle{plain} %              Remove in Eng. 


\begin{document} 


%\thispagestyle{plain}%          Set in Eng. 


\maketitle 


\label{begin} 


 


 


 


Под {\it базисом} понимаем конечную полную систему булевых функций. 


 


Булева функция называется {\it бесповторной} в базисе $B$, если ее можно 


представить в этом базисе термом (формулой), в котором каждая 


переменная встречается не более одного раза. В противном случае она 


называется повторной в $B$. 


 


Булева функция называется {\it слабоповторной} в базисе $B$, если ее любая 


остаточная подфункция является бесповторной, а сама она повторна в базисе 


$B$. 


 


Под {\it сложностью} $L(\Phi)$ терма $\Phi$ понимаем число всех вхождений 


переменных 


в $\Phi$. {\it Сложностью} $L_B(f)$ булевой функции $f$ в базисе $B$ 


называется 


наименьшее значение $L(\Phi)$ при условии, что терм $\Phi$ в базисе 


$B$ представляет функцию $f$. 


 


При сравнении базисов по сложности представлений булевых функций 


термами на множестве всех базисов вводится частичный порядок: 


$B_1\le B_2$, если существует число $c$ такое, что $L_{B_1}(f)\le c 


L_{B_2}(f)$ 


для любой булевой функции $f$, говорят, $B_1$ {\it предшествует} $B_2$. 


Если 


$B_1\le B_2$ и $B_2\le B_1$, то базисы $B_1$ и $B_2$ называются {\it 


эквивалентными}. 


Если $B_1\le B_2$ и $B_2\not\le B_1$, то пишем $B_1<B_2$ и говорим, что 


$B_1$ 


{\it строго предшествует} $B_2$. Также говорим, что $B_1$ {\it 


непосредственно 


предшествует} $B_2$, если $B_1<B_2$ и не существует базиса $B$ такого, что 


$B_1<B<B_2$. 


 


Таким образом, множество базисов разбито на классы эквивалентности. 


В [1] доказано, что в каждом классе базисов можно указать канонический 


вид класса, причем если базис $B$ непосредственно предшествует 


базису $B'$, то канонический вид $B$ содержит на одну функцию больше, 


чем канонический вид $B'$, а эта функция является слабоповторной в 


базисе, каноническом для $B'$. 


 


О.Б.\,Лупановым замечено (см. [2]), что базис $B_0=\{\vee,\cdot,-,0,1\}$ 


является наибольшим элементом при введенном 


порядке, т.\,е. класс базисов, эквивалентных базису $B_0$, самый ``плохой'' 


по сложности реализации булевых функций термами. Эти базисы 


назовем базисами {\it нулевого яруса}. Отметим, что базис $B_0$ 


канонический 


для своего класса. 


 


Базисами $k$-{\it го яруса} ($k>0$) называются все базисы, непосредственно 


предшествующие всем базисам $(k-1)$-го яруса. 


 


 


Введем обозначение $x^\sigma=\begin{cases} 


x, & \text{если $\sigma=1$};\\ 


\ov x, & \text{если $\sigma=0$}. 


\end{cases}$ 


 


\noindent{}Функции $f$ и $g$ называются однотипными, 


если $f(x_1,\dotsc,x_n)=g(x^{\sigma_1}_{i_1},\dotsc,x^{\sigma_n}_{i_n})$, 


где $(i_1,\dotsc,i_n)$ --- 


произвольная перестановка чисел от 1 до $n$. 


 


Функции $f$ и $g$ называются {\it обобщенно однотипными}, если $f$ 


однотипна к $g$ или $\ov g$. Очевидно, что на множестве булевых функций 


отношение 


обобщенной однотипности является отношением эквивалентности. 


 


В работе [3] получено описание всех обобщенных типов функций, 


слабоповторных в базисе $B_0$, а как следствие --- канонических базисов 


первого яруса. Заметим, что при добавлении к базису $B_0$ разных по 


типу слабоповторных функций получаются базисы из разных классов 


эквивалентности. Поэтому уже в первом ярусе имеется счетное число различных 


классов эквивалентности базисов. 


 


Естественно, возник вопрос изучения базисов второго яруса. Часть 


этих базисов удалось описать в [4]--[7]. 


 


К настоящему моменту автором найдены все слабоповторные функции 


в канонических базисах первого яруса, тем самым, описание базисов 


второго яруса завершено. 


 


Ниже следует один из результатов, полученных автором по данной 


проблеме. 


 


\begin{thmnonum} 


Система булевых функций 


\begin{alignat*}{2} 


& x_1(x_2\vee x_3x_4)\vee x_5(x_3\vee x_2x_4), &\quad & \ \\ 


% 


&x_1(x_2\vee\dotsb\vee x_n)\vee x_2\cdot\dotsc\cdot x_n, 


&\quad & n\ge3, \\ 


% 


&x_1(x_2\vee x_3\cdot\dotsc\cdot x_n)\vee x_2\cdot\ov x_3\cdot 


\dotsc \cdot\ov x_n, &\quad & n\ge 3, \\ 


% 


&x_1\cdot\dotsc\cdot x_n\vee \ov x_1\cdot\dotsc\cdot\ov x_n, 


&\quad & n\ge 2, \\ 


% 


&\ov x_1g(x_2,\dotsc,x_5)\vee x_1g(x_3,x_2,x_5,x_4), &\quad & \ \\ 


% 


&\ov x_1g(x_2,\dotsc,x_5)\vee x_1x_3x_5(x_2\vee x_4), &\quad & \ \\ 


% 


&\ov x_1g(x_2,\dotsc,x_5)\vee x_1(x_2x_3\vee x_4x_5), &\quad & \ \\ 


% 


&\ov x_1g(x_2,\dotsc,x_5)\vee x_1x_3(x_2\vee x_4x_5), &\quad & \ \\ 


% 


&\ov x_1x_2g(x_3,\dotsc,x_6)\vee x_1g 


(x_2x_3,x_4,x_5,x_6) &\quad & \ 


\end{alignat*} 


является полной системой представителей классов эквивалентности 


по отношению обобщенной однотипности для функций, слабоповторных 


в базисе $B=B_0\cup\{g\}$, где $g(x_1,\dotsc,x_4)=x_1(x_2\vee x_3)\vee 


x_3x_4$. 


\end{thmnonum} 


 


 


\begin{thebibliography}{9} 


 


\bibitem{1} 


Черухин Д.Ю. {\em Алгоритмический критерий сравнения булевых базисов} 


// Матем. вопр. кибернетики. -- 1999. -- \No\,8. -- С.\,77--122. 


\bibitem{2} 


Субботовская Б.А. {\em О сравнении базисов при реализации функций 


алгебры логики формулами} // ДАН СССР. -- 1963. -- Т.\,149. -- \No\,4. 


-- С.\,784--787. 


\bibitem{3} 


Стеценко В.А. {\em О предплохих базисах в $P_2$} // Матем. вопр. 


кибернетики. -- 1992. -- \No\,4. -- С.\,139--177. 


\bibitem{4} 


Перязев Н.А. {\em Слабоповторные булевы функции в бинарном базисе} 


// Дискретн. матем. и информатика. -- Иркутск. ун-т, 1998. -- Вып.\,4. -- 


12 с. 


\bibitem{5} 


Кириченко К.Д. {\em Слабоповторные булевы функции в некоторых 


предэлементарных базисах} // 


Дискретн. матем. и информатика. -- Иркутск. ун-т, 2000. -- Вып.\,13. -- 


60 с. 


\bibitem{6} 


Шаранхаев И.К. {\em Слабоповторные булевы функции в одном предэлементарном 


базисе} // Методы оптимизации и их прилож. 


Тр. 12-й Байкал. межд. конф. -- 2001. -- Т.\,5. -- С.\,151--155. 


\bibitem{7} 


Шаранхаев И.К. {\em О слабоповторных булевых функциях в одном базисе} 


// Дискретн. анализ и исслед. операций. Матер. 


конф. -- Новосибирск, 2002. -- С.\,139. 


 


\end{thebibliography} 


 


\vskip 2cm 


 


\post{\it Бурятский государственный}{\it Поступила} 


\post{\it университет}{21.03.2003} 


 


\label{end} 


\end{document} 


